
\subsection{Literal comparison with the earlier finite source and its cocycle}

The earlier source (TC1) writes the spectral polynomial coordinate as
$t$ and the same residue representative as $R_Z(u)$; the coordinate
isomorphism sends $t$ to $s$, with $h_Z(t)$ sent to $h(s)$, every
power $(t-\rho)^r$ sent to $(s-\rho)^r$, and the full unit unchanged.
Consequently $R_Z(u)$ is $P_h(u)$ under this specified isomorphism.
For the actual source functions, not only their quotient classes,
\[
R_{\mathrm{ref}}=s_h:E_h\longrightarrow\mathscr B_X,
\qquad X=U\text{ or }\Omega\text{ as the packet requires}.
\tag{MCF42}
\]
Indeed applying $h(D)$ to either representative gives
$\Theta(P_h(u)(D)\phi_*)$.  The difference lies in $\mathscr B$
and is killed by the injective $h(D)$, hence vanishes.  This proves
the asserted equality with the original units and domains.

Put $t=\log a$, retaining its orientation when $a<1$, and define
\[
c_{t,h}(u)=\int_0^t\mathcal R_{e^{t-r}}\phi_*\,
                  \ell_h(e^{rA_h}u)\,dr:E_h\longrightarrow V.
\tag{MCF43}
\]
The integral converges in the original Schwartz topology.  In fact
on the compact interval joining $0$ and $t$, every derivative of
$\mathcal R_{e^{t-r}}\phi_*$ is a continuously varying rescaled
derivative of the same Schwartz function, with the scaling parameter
in a compact subset of $(0,\infty)$.  For every pair $N,k$, its
seminorm $\sup_x(1+|x|)^N|\partial_x^k(\cdot)|$ is uniformly
bounded there.  The finite matrix $e^{rA_h}$ is smooth and bounded
there too.  Riemann sums therefore converge in every seminorm in
the complete Schwartz space; parity and the two zero moments pass
to the limit by continuity, so the limit is in $V$.

Differentiate the actual $\mathscr B$-valued function
$\mathcal R_{e^{t-r}}s_he^{rA_h}u$.  Its derivative is
$-\mathcal R_{e^{t-r}}(Ds_h-s_hA_h)e^{rA_h}u$.
The same seminorm argument on $\mathscr B$, followed by (MCF5),
allows the fundamental theorem of calculus.  It gives
\[
\Theta c_{\log a,h}=\mathcal R_as_h-s_hA_{a,h},\qquad
\Psi_{a,h}=c_{\log a,h}.
\tag{MCF44}
\]
The second equality follows by comparing with (MCF24) and canceling
the injective map $\Theta:V\to\mathscr B$.  Thus the previous
oriented integral and the Euler-inverse primitive are the same map.
Their packet and coefficient-face naturality are literal, by
(MCF27) and coordinatewise integration.  The joint primitive is
$(c/2,-\widehat c/2)$; the difference between the two single-leg
primitives $(c,0)$ and $(0,-\widehat c)$ is the full cycle
$(c,\widehat c)$.  No diagonal is discarded by this comparison.

The earlier dual two-leg complex is obtained by applying algebraic
$\operatorname{Hom}_{\mathbb C}(-,\mathbb C)$, so the map induced
by $C_U\to C_+$ is restriction $\mathscr B^\vee\to\mathscr B_U^\vee$
in degree $-1$ and identity on both $V^\vee$ source coordinates in
degree zero.  Its differential, in the same order of legs, is
\[
\ell\longmapsto(\ell\Theta,-\ell\Theta\mathcal F),
\qquad \mathcal F\phi=\widehat\phi.
\tag{MCF45}
\]
Because $\Theta V\subseteq\mathscr B_U$, restriction followed by
this differential equals the differential followed by the identity.
Thus this is an actual cochain map.  It induces restriction
$Q^\vee\to(T_UQ)^\vee$ on degree $-1$ cohomology.  At a joint
source the two diagonal parameters $v\mapsto(v,\mathcal Fv)$ and
$w\mapsto(\mathcal Fw,w)$ agree under $w=\mathcal Fv$. Define
$\mathcal F^\vee:V^\vee\to V^\vee$ by
$\mathcal F^\vee(\ell)=\ell\circ\mathcal F$; the maps
$\pi_v,\pi_w:V^\vee\oplus V^\vee\to V^\vee$ are
\[
\pi_v(\alpha,\beta)=\alpha+\beta\mathcal F,
\quad \pi_w(\alpha,\beta)=\alpha\mathcal F+\beta,
\quad \pi_v=\mathcal F^\vee\circ\pi_w.
\tag{MCF46}
\]
These identities follow by evaluation on the displayed diagonals
and $\mathcal F^2=1$.  The original actions intertwine because
$\mathcal F\mathcal R_a=a\mathcal R_{1/a}\mathcal F$.
For a proper source $W$, dualizing $W\hookrightarrow V$ gives
restriction of diagonal functionals to $W$.  For coefficient faces,
dualizing zero insertion gives coordinate restriction.  These are
algebraic dual statements; no continuous splitting for arbitrary
nonclosed $W$ is asserted.
